{\centering\tiny\renewcommand{\arraystretch}{0.92}
\begin{tabularx}{\textwidth}{@{}>{\raggedright\arraybackslash}p{2.6cm}Y@{}}
\toprule
family & rules (id: name [status]) \\
\midrule
\textsc{pos} (8) position, scope, normalisation & \textbf{01} Prāsa position (2nd akshara of every pāda) [C]; \textbf{02} Meter eligibility [C]; \textbf{03} Uniformity across meters [I]; \textbf{04} Stanza scope [C]; \textbf{05} Sandhi surface form is what is judged [C]; \textbf{06} Saṃśleṣa [M]; \textbf{07} Druta-sandhi alternative reading [C]; \textbf{08} Dantya \te{ౘ/ౙ} normalisation [C] \\
\textsc{sama} (12) samaprāsa: identity and neutral features & \textbf{01} Samaprāsa [C]; \textbf{02} Vowel neutrality [C]; \textbf{03} \te{ఋ}-tva sahita [C]; \textbf{04} \te{ఌ}-kāra sahita [C]; \textbf{05} Bindu-sahita [C]; \textbf{06} Visarga-sahita [C]; \textbf{07} Drutam following the prāsa akshara is neutral [C]; \textbf{08} Pollu \te{ల్} (or any dead consonant) following the prāsa akshara is neutral [C]; \textbf{09} Dantya and tālavya \te{చ/జ} are savarṇa [C]; \textbf{10} Weight of the prāsa akshara itself is free [I]; \textbf{11} Laghu-yakāra prāsa [Cs]; \textbf{12} Sandhigata prāsa [Cs] \\
\textsc{samyukta} (8) conjuncts & \textbf{01} Conjunct presence must agree (saṃyuktākṣara prāsamu) [M]; \textbf{02} Constituent consonants and their order must be identical (kramamu) [M]; \textbf{03} Dvitvākṣara prāsa [M]; \textbf{04} Three-consonant clusters keep all three in order [M]; \textbf{05} \te{ఋ} inside a conjunct is a vowel [C]; \textbf{06} Pūrṇabindu-pūrvaka saṃyuktākṣara prāsa [Cs]; \textbf{07} Drutam-fused conjunct prāsa (\te{న్}X) [Cs]; \textbf{08} Pollu-\te{ల్}-fused conjunct prāsa (\te{ల్}X) [Cs] \\
\textsc{purvaka} (6) signs before the prāsa akṣara & \textbf{01} Pūrṇabindu prāsa [M]; \textbf{02} Visarga prāsa [M]; \textbf{03} Ardhabindu prāsa [Cs]; \textbf{04A} Khaṇḍākhaṇḍa prāsa within Ananta's restriction (\te{ప/ట}) [Cs]; \textbf{04B} Khaṇḍākhaṇḍa prāsa within Appakavi's restriction (paruṣa after a long vowel) [Cs]; \textbf{04C} Khaṇḍākhaṇḍa prāsa in its wide classical form [R] \\
\textsc{purvakshara} (4) weight of the pre-prāsa akṣara & \textbf{01} Pre-prāsa akshara weight uniformity [M]; \textbf{02} Metrical weight, not vowel length [C]; \textbf{03} Kandamu gana corollary [I]; \textbf{04} Light repha clusters [Cs] \\
\textsc{maitri} (15) accepted or deprecated equivalences & \textbf{SA-SHA} \te{స} $\sim$ \te{శ} sibilant compatibility [R]; \textbf{THA-DHA} \te{థ} $\sim$ \te{ధ} svavargaja prāsa [R]; \textbf{NA-NNA} \te{న} $\sim$ \te{ణ} ubhaya prāsa (1) [R]; \textbf{LA-LLA} \te{ల} $\sim$ \te{ళ} abhēda prāsa [R]; \textbf{RU-RA} \te{ఋ} prāsa [R]; \textbf{RU-HALSAMYUTA} Hal-saṃyuta \te{ఋ} prāsa [R]; \textbf{VIKALPA} Vikalpa prāsa [R]; \textbf{BINDU-SAMSLESHA} Bindu $\sim$ saṃśleṣa spelling of one nasal + stop (\te{ంద} $\sim$ \te{న్ద}) [R]; \textbf{RA-RRA-ORTHO} \te{ర} $\sim$ \te{ఱ} orthographic relaxation (rēpha $\sim$ śakaṭarēpha) [O]; \textbf{SA-SSA} \te{స} $\sim$ \te{ష} ubhaya prāsa (2) [D]; \textbf{LA-DA} \te{ల/ళ} $\sim$ \te{డ} [D]; \textbf{DA-DHA} \te{ద} $\sim$ \te{ధ} [D]; \textbf{SAMYUKTASAMYUKTA} Saṃyuktāsaṃyukta prāsa [D]; \textbf{ANUNASIKA} Anunāsika prāsa [D]; \textbf{MBA-MMA} \te{ంబ} $\sim$ \te{మ్మ} prāsamaitri prāsa (deprecated) [D] \\
\textsc{rephayuta} (1) repha conjuncts & \textbf{KRARAKOMMU} Krārakommu vs vaṭruvasuḍi may not rhyme [F] \\
\textsc{vaira} (7) forbidden pairs & \textbf{RA-RRA} \te{ర} $\neq$ \te{ఱ} (prāsavairamu) [F]; \textbf{TA-DA} \te{ట} $\neq$ \te{డ} [F]; \textbf{DA-DDA} \te{ద} $\neq$ \te{డ} [F]; \textbf{DDA-DDHA} \te{డ} $\neq$ \te{ఢ} [F]; \textbf{RA-LA} \te{ర} $\neq$ \te{ల} (niṣiddhamu) [F]; \textbf{SVAVARGAJA-EXT} Svavargaja licence does not extend beyond \te{థ-ధ} [F]; \textbf{GENERIC} No prāsamaitri between these consonants (asamaprāsa) [F] \\
\textsc{defect} (3) named defects & \textbf{PURNARDHABINDU} Pūrṇārdhabindu prāsa [F]; \textbf{ADHIKA} Adhika prāsa [F]; \textbf{SANTA} Śānta prāsa [F] \\
\bottomrule
\end{tabularx}
\captionof{table}{The prāsa rules (\texttt{prasa\_rules.yaml}), distilled from our notes on the classical treatises, among them the \textit{Appakavīyamu} \citep{appakavi1656} (64 rules). Status: C canonical, Cs canonical subtype, M mandatory (all profiles), R accepted relaxation (relaxed and historical profiles), O orthographic relaxation, D deprecated (historical only), F forbidden (a named failure), I informational.}
\label{tab:prasarules}\par}
\medskip
